\section{Introduction}

The conjugate gradient method (CG) of Hestenes and Stiefel~\cite{hestenes1952} is the
canonical Krylov solver for a symmetric positive definite (SPD) linear system $A x = b$;
see, e.g., Golub and Van Loan~\cite{golub2013} or Greenbaum~\cite{greenbaum1997}. It builds
its iterates in the Krylov subspace $\mathcal{K}_k(A, b)$ and converges at the
$O(\sqrt\kappa)$ rate, where $\kappa = \kappa(A)$ is the spectral condition number
(Proposition~\ref{prop:cg}). What CG does \emph{not} do is respect an inequality: it solves
equations, not constrained programs. This note is about the gap between the two -- how to
make CG deliver the solution of the bound-constrained quadratic
\begin{equation}\label{eq:bqp}
  \min_{x \in \mathbb{R}^n}\; f(x) = \tfrac{1}{2}\,x^\top A x - b^\top x
  \quad\text{subject to}\quad x \geq 0,
  \qquad A \succ 0,
\end{equation}
without giving up its per-iteration cost or its Krylov convergence rate.

Problem~\eqref{eq:bqp} is the strictly convex non-negative quadratic program. When
$A = M^\top M$ is the Gram matrix of $M \in \mathbb{R}^{m\times n}$ and $b = M^\top d$, its
minimiser is the solution of the non-negative least-squares (NNLS) problem
$\min_{x\geq 0}\|Mx - d\|_2^2$, the setting of Lawson and Hanson~\cite{lawson1974}. Without
the constraint, the minimiser is simply $x = A^{-1}b$, which is exactly what CG computes; the
constraint $x \geq 0$ is what stands between that clean solve and the answer. We also treat
the \emph{equality-augmented} variant, in which a single linear normalisation
$\mathbf{1}^\top x = \beta$ is imposed alongside $x \geq 0$, because it arises whenever the
solution is required to lie on an affine slice of the non-negative orthant and because its
multiplier can be eliminated analytically (Section~\ref{sec:reduction}).

This note studies the constraint-handling layer directly. Its contribution is a
\emph{primal-dual active-set} outer loop that wraps CG: it solves the unconstrained system
over a working set of \emph{free} variables, releases any variable that returns negative
(the primal step), re-admits any bound variable whose reduced gradient is negative (the dual
step), and restarts CG on the modified free set each time. We make three points about this
construction. First, the variable toggles are precisely the \emph{principal pivots} of the
linear complementarity problem (LCP) $(A,\,{-b})$, and $A \succ 0$ makes it a $P$-matrix, so
every principal submatrix is invertible and each pivot is well
defined~\cite{murty1988,judice1994,kim2011}. Second, guarding a fast block-pivot path with a
least-index Bland-style fallback yields \emph{unconditional} finite termination at the unique
global minimiser, with no non-degeneracy hypothesis (Theorem~\ref{thm:termination}). Third,
integrating a Krylov inner solver with a changing free set -- restarting CG on each reduced
system and warm-starting it across a parametric family of problems -- is what turns the
finite-termination guarantee into a fast solver in practice.

The inner solve preserves two properties that matter for speed. It is \emph{matrix-free}:
when $A = M^\top M$, the action $v \mapsto A_{\mathcal{F}}v = M_{\mathcal{F}}^\top(M_{\mathcal{F}}v)$
on the free set $\mathcal{F}$ is evaluated as two products with $M_{\mathcal{F}}$, so $A$ is
never assembled and no $O(n^2)$ storage is incurred. And its iteration count is governed by
$\kappa$, so any preconditioner that compresses the spectrum accelerates it. A ridge/Tikhonov
splitting $A \mapsto A_\alpha = (1-\alpha)A + \alpha R^\top R$ with $R^\top R \succ 0$ does
both: it lowers $\kappa$ (Proposition~\ref{prop:regcond}) and, by keeping every principal
submatrix strictly positive definite, secures the $P$-matrix property the finite-termination
proof relies on (Section~\ref{sec:conditioning}).

For contrast we analyse a loop-free comparator that enforces the constraint at every step: a
\emph{projected-gradient} method that takes a gradient step and projects onto the feasible set
-- the non-negative orthant for~\eqref{eq:bqp}, or the affine simplex slice
$\Delta_\beta = \{x : \mathbf{1}^\top x = \beta,\, x \geq 0\}$ for the equality-augmented
variant, the latter computed in $O(n\log n)$ by a sort-and-threshold
pass~\cite{duchi2008,condat2016}. It is single-stage and simple, but it is not a Krylov
method: it does not orthogonalise residuals against prior search directions and converges at
the $O(\kappa)$ rate, a factor $\sqrt\kappa$ slower than the CG inner loop
(Section~\ref{sec:nonneg}).

The two ingredients are individually classical -- active-set and block-principal-pivoting
methods for the non-negative and complementarity problems~\cite{lawson1974,judice1994,kim2011,%
murty1988,nocedal2006}, and CG with preconditioning for SPD systems~\cite{hestenes1952,%
golub2013,greenbaum1997,benzi2005}. What this note isolates is their combination: a
matrix-free Krylov inner solver driven through a changing active set with an unconditional
finite-termination guarantee, and warm-started across a parametric sweep. The algorithm is
the computational kernel of the companion papers~\cite{schmelzer2026matrixfree,%
schmelzer2026rmt}, which instantiate it on large structured problems; here it is developed and
analysed on its own, independent of any application.

Section~\ref{sec:problem} states the problem and its KKT/complementarity structure;
Section~\ref{sec:reduction} reduces it to a sequence of unconstrained SPD solves, records the
matrix-free CG operator and its convergence rate, and eliminates the equality multiplier
analytically; Section~\ref{sec:nonneg} develops the active-set loop and the projected-gradient
comparator and proves finite termination; Section~\ref{sec:conditioning} treats regularisation
and preconditioning; and Section~\ref{sec:results} records the complexity of each method and
discusses its behaviour on synthetic SPD test problems.
